My Hermite work taught me to evaluate an action representation under a bounded offline protocol, while CR5 deployment showed how an observation-command temporal mismatch under non-blocking execution produced unstable corrections. I can investigate these layers separately, but my capability gap is broader grounding in artificial-intelligence foundations and robust evaluation that connects algorithms to deployed behavior. CUHK's MSc in Artificial Intelligence is well matched to this need. If offered in my intake, AIMS5701 Fundamentals in Artificial Intelligence would consolidate the conceptual basis beneath my current implementation experience, and AIST5020 Trustworthy Artificial Intelligence would sharpen how I define evidence, failure modes, and reliability claims.

ROSE5760 Robot Learning would let me study action representation within the wider problem of learning behavior, rather than treating token fidelity as a complete policy result. Subject to approval and availability, AIMS5790 Artificial Intelligence Project I would provide a setting to test a focused question across data, representation, and execution, informed by both my Hermite branch and WebSocket-served CR5 policy loop. I would contribute a habit of separating codec evidence from rollout and hardware evidence, and use the programme to develop AI systems whose evaluation remains meaningful when models interact with physical agents in robotics research and development.
